% ============================================================
%  styles/commands.tex
%  Unified notation for Semantic Infrastructure monograph
% ============================================================

% ── RSVP field triple ───────────────────────────────────────
\newcommand{\rsvp}{(\Phi,\,\mathbf{v},\,S)}        % full triple
\newcommand{\rsvpfull}{\bigl(\Phi,\,\mathbf{v},\,S\bigr)}
\newcommand{\Phifield}{\Phi}                         % scalar field
\newcommand{\vfield}{\mathbf{v}}                     % vector field
\newcommand{\Sfield}{S}                              % entropy field

% ── Accessibility & admissibility ───────────────────────────
\newcommand{\Acc}{\mathcal{A}}                       % accessibility set
\newcommand{\Adm}{\mathbf{Adm}}                      % admissibility operator
\newcommand{\proj}{\pi}                              % projection
\newcommand{\Proj}{\Pi}                              % projection operator
\newcommand{\fiber}[2]{#1^{-1}(#2)}                 % fiber over a point
\newcommand{\constraint}{\mathcal{C}}                % constraint space
\newcommand{\admissible}[1]{\mathrm{Adm}(#1)}

% ── Spherepop notation ──────────────────────────────────────
\newcommand{\pop}{\mathbin{\bullet}}                 % pop operator
\newcommand{\sphere}[1]{\langle #1 \rangle}          % sphere containment
\newcommand{\nest}[2]{\langle #1 \mid #2 \rangle}    % nested sphere
\newcommand{\eval}[1]{\llbracket #1 \rrbracket}      % evaluation brackets
\newcommand{\collapse}[1]{\downarrow\!#1}            % collapse operator
\newcommand{\refuse}{\bot_{\mathrm{sp}}}             % refusal

% ── CLIO ────────────────────────────────────────────────────
\newcommand{\clio}{\textsf{CLIO}}
\newcommand{\tartan}{\textsf{TARTAN}}
\newcommand{\rsvpname}{\textsf{RSVP}}

% ── Categorical notation ────────────────────────────────────
\newcommand{\cat}[1]{\mathbf{#1}}                    % category
\newcommand{\ob}[1]{\mathrm{Ob}(#1)}                 % objects
\newcommand{\Hom}[3]{\mathrm{Hom}_{#1}(#2,#3)}      % hom-set
\newcommand{\Nat}{\mathrm{Nat}}                      % natural transformations
\newcommand{\fun}[2]{#1 \Rightarrow #2}              % functor arrow
\newcommand{\sheaf}{\mathcal{F}}                     % sheaf
\newcommand{\stalk}[2]{\mathcal{F}_{#2}}             % stalk at point

% ── Geometry & manifolds ────────────────────────────────────
\newcommand{\manifold}[1]{\mathcal{M}_{#1}}
\newcommand{\tangent}[1]{T_{#1}}
\newcommand{\cotangent}[1]{T^*_{#1}}
\newcommand{\gradient}{\nabla}
\newcommand{\divg}{\nabla \cdot}
\newcommand{\curl}{\nabla \times}
\newcommand{\laplacian}{\Delta}
\newcommand{\dalembertian}{\square}

% ── Entropy & thermodynamics ─────────────────────────────────
\newcommand{\entropy}{S}
\newcommand{\dS}{\mathrm{d}S}
\newcommand{\lamphron}{\lambda_{\!+}}               % lamphron
\newcommand{\lamphrodyne}{\lambda_{\!-}}            % lamphrodyne

% ── Gesture & trajectory ────────────────────────────────────
\newcommand{\gesture}[1]{\gamma_{#1}}                % gesture trajectory
\newcommand{\gestureidx}{\mathcal{G}}           % gesture subscript (no arg)
\newcommand{\gesturesp}{\mathcal{G}}                 % gesture space
\newcommand{\motorclass}[1]{[#1]_{\sim}}             % motor equiv class
\newcommand{\canonical}{\hat{\gamma}}                % canonical gesture

% ── Action history ──────────────────────────────────────────
\newcommand{\action}{\mathcal{S}}                    % action functional
\newcommand{\history}{\mathcal{H}}                   % action history
\newcommand{\commitment}{\kappa}                     % commitment variable
\newcommand{\freedom}{\phi}                          % remaining freedom

% ── Text shortcuts ──────────────────────────────────────────
\newcommand{\ie}{\textit{i.e.},\ }
\newcommand{\eg}{\textit{e.g.},\ }
\newcommand{\cf}{\textit{cf.}\ }
\newcommand{\viz}{\textit{viz.}\ }

% ── Margin annotation ───────────────────────────────────────
\newcommand{\sidenote}[1]{\marginnote{\footnotesize\color{RSVPmid}\itshape #1}}
\newcommand{\sidedef}[1]{\marginnote{\footnotesize\color{RSVPaccent}\bfseries #1}}

% ── Chapter epigraph ────────────────────────────────────────
\newcommand{\chapterepigraph}[2]{%
  \begin{flushright}
    \begin{minipage}{0.55\linewidth}
      \small\itshape #1\\[2pt]
      \upshape\footnotesize\color{RSVPmid}--- #2
    \end{minipage}
  \end{flushright}
  \vspace{1em}
}

% ── Named theorems convenience wrappers ─────────────────────
% Usage: \begin{namedthm}{Accessibility Relaxation Theorem}
\newenvironment{namedthm}[1]{%
  \begin{tcolorbox}[
    enhanced, breakable,
    colback=RSVPlight,
    colframe=RSVPdeep,
    title={\bfseries\sffamily #1},
    coltitle=white,
    attach boxed title to top left={yshift=-2mm,xshift=6mm},
    boxed title style={colback=RSVPdeep, sharp corners},
    arc=4pt, boxrule=1pt,
    top=4mm, left=4mm, right=4mm, bottom=4mm,
  ]
}{%
  \end{tcolorbox}
}
